% GENERATED FILE. Edit canonical lessons, then run npm run generate:books.
% canonical-source-hash: fnv1a64:5723648412cdc2cd
% canonical-lessons: HI-C246-loha, HI-C246-cadi, HI-C246-kagaz, HI-C246-plastik, HI-C246-rui

\chapter{Iron, Silver, Paper, Plastic, Cotton wool}
\label{ch:hi-iron-silver-paper-plastic-cotton-wool}
\hlchaptermodality{246}

\begin{chapteropening}
\textbf{By the end of this chapter:} \emph{I can say iron, silver, paper, plastic and cotton wool.}

Complete the last lesson of chapter 246: I can say iron, silver, paper, plastic and cotton wool.
\end{chapteropening}

\section[\texorpdfstring{\hi{लोहा} (lohā)}{लोहा (lohā)}]{\hi{लोहा} (lohā) — iron}

\label{lesson:HI-C246-loha}



\begin{quote}
Before the new one: say the Hindi for a dozen, then the Hindi for a pair.
\end{quote}

\subsection*{You'll want to know: \hi{लोहा}}
\textbf{\hi{लोहा}} — \emph{lohā} — \textquotedblleft{}iron\textquotedblright{}.

\begin{grammarlens}[title={What you've built}]
One more word for this chapter.
\end{grammarlens}

\subsection*{Your turn}
\begin{itemize}
\raggedright
  \item \emph{Say it:} \emph{lohā}
  \item \emph{Say it:} \emph{lohā}, once more
  \item \emph{Say it:} say \emph{lohā}
  \item \emph{Recall:} say the Hindi for a dozen, then the Hindi for a pair, then say \emph{lohā} again
\end{itemize}

\begin{tcolorbox}[breakable,colback=teal!4,colframe=teal!35!black,title={Before you move on}]
What does \hi{लोहा} mean? (\textquotedblleft{}Iron\textquotedblright{}.) Say it once more.
\end{tcolorbox}
\section[\texorpdfstring{\hi{चाँदी} (cā̃dī)}{चाँदी (cā̃dī)}]{\hi{चाँदी} (cā̃dī) — silver}

\label{lesson:HI-C246-cadi}



\begin{quote}
Before the new one: say the Hindi for a pair, then the Hindi for iron.
\end{quote}

\subsection*{You'll want to know: \hi{चाँदी}}
\textbf{\hi{चाँदी}} — \emph{cā̃dī} — \textquotedblleft{}silver\textquotedblright{}.

\begin{grammarlens}[title={What you've built}]
One more word for this chapter.
\end{grammarlens}

\subsection*{Your turn}
\begin{itemize}
\raggedright
  \item \emph{Say it:} \emph{cā̃dī}
  \item \emph{Say it:} \emph{cā̃dī}, once more
  \item \emph{Say it:} say \emph{cā̃dī}
  \item \emph{Recall:} say the Hindi for a pair, then the Hindi for iron, then say \emph{cā̃dī} again
\end{itemize}

\begin{tcolorbox}[breakable,colback=teal!4,colframe=teal!35!black,title={Before you move on}]
What does \hi{चाँदी} mean? (\textquotedblleft{}Silver\textquotedblright{}.) Say it once more.
\end{tcolorbox}
\section[\texorpdfstring{\hi{काग़ज़} (kāgaz)}{काग़ज़ (kāgaz)}]{\hi{काग़ज़} (kāgaz) — paper}

\label{lesson:HI-C246-kagaz}



\begin{quote}
Before the new one: say the Hindi for iron, then the Hindi for silver.
\end{quote}

\subsection*{You'll want to know: \hi{काग़ज़}}
\textbf{\hi{काग़ज़}} — \emph{kāgaz} — \textquotedblleft{}paper\textquotedblright{}.

\begin{grammarlens}[title={What you've built}]
One more word for this chapter.
\end{grammarlens}

\subsection*{Your turn}
\begin{itemize}
\raggedright
  \item \emph{Say it:} \emph{kāgaz}
  \item \emph{Say it:} \emph{kāgaz}, once more
  \item \emph{Say it:} say \emph{kāgaz}
  \item \emph{Recall:} say the Hindi for iron, then the Hindi for silver, then say \emph{kāgaz} again
\end{itemize}

\begin{tcolorbox}[breakable,colback=teal!4,colframe=teal!35!black,title={Before you move on}]
What does \hi{काग़ज़} mean? (\textquotedblleft{}Paper\textquotedblright{}.) Say it once more.
\end{tcolorbox}
\section[\texorpdfstring{\hi{प्लास्टिक} (plāsṭik)}{प्लास्टिक (plāsṭik)}]{\hi{प्लास्टिक} (plāsṭik) — plastic}

\label{lesson:HI-C246-plastik}



\begin{quote}
Before the new one: say the Hindi for silver, then the Hindi for paper.
\end{quote}

\subsection*{You'll want to know: \hi{प्लास्टिक}}
\textbf{\hi{प्लास्टिक}} — \emph{plāsṭik} — \textquotedblleft{}plastic\textquotedblright{}.

\begin{grammarlens}[title={What you've built}]
One more word for this chapter.
\end{grammarlens}

\subsection*{Your turn}
\begin{itemize}
\raggedright
  \item \emph{Say it:} \emph{plāsṭik}
  \item \emph{Say it:} \emph{plāsṭik}, once more
  \item \emph{Say it:} say \emph{plāsṭik}
  \item \emph{Recall:} say the Hindi for silver, then the Hindi for paper, then say \emph{plāsṭik} again
\end{itemize}

\begin{tcolorbox}[breakable,colback=teal!4,colframe=teal!35!black,title={Before you move on}]
What does \hi{प्लास्टिक} mean? (\textquotedblleft{}Plastic\textquotedblright{}.) Say it once more.
\end{tcolorbox}
\section[\texorpdfstring{\hi{रुई} (rūī)}{रुई (rūī)}]{\hi{रुई} (rūī) — cotton wool}

\label{lesson:HI-C246-rui}



\begin{quote}
Before the new one: say the Hindi for paper, then the Hindi for plastic.
\end{quote}

\subsection*{You'll want to know: \hi{रुई}}
\textbf{\hi{रुई}} — \emph{rūī} — \textquotedblleft{}cotton wool\textquotedblright{}.

\begin{grammarlens}[title={What you've built}]
One more word for this chapter.
\end{grammarlens}

\subsection*{Your turn}
\begin{itemize}
\raggedright
  \item \emph{Say it:} \emph{rūī}
  \item \emph{Say it:} \emph{rūī}, once more
  \item \emph{Say it:} say \emph{rūī}
  \item \emph{Recall:} say the Hindi for paper, then the Hindi for plastic, then say \emph{rūī} again
\end{itemize}

\begin{tcolorbox}[breakable,colback=teal!4,colframe=teal!35!black,title={Before you move on}]
What does \hi{रुई} mean? (\textquotedblleft{}Cotton wool\textquotedblright{}.) Say it once more.
\end{tcolorbox}
